%!TeX root=csp-audit.tex
\section{The legislated readings}\label{sec:readings}

Ten places where \cite{ZhukSimplified} admits more than one reading. The choice
is invisible in prose and forced in a formalization, and in six of the ten it is
not really a choice: one reading makes some statement of the source true or
some proof valid and the other does not, so the ``convention'' is a
reconstruction. Two of the ten are not conventions at all --- one is a theorem
and one is a missing lemma --- and only two are free.

\begin{center}\small
\begin{tabular}{@{}llp{0.46\textwidth}@{}}
\hline
\S & \textbf{Item} & \textbf{Status} \\
\hline
\ref{sec:leg-Zp} & the algebra $\mathbf Z_{p}$ & \textbf{a theorem}, proved \\
\ref{sec:leg-cover} & $\cover\sigma$ is a tolerance & decided by internal evidence \\
\ref{sec:leg-empty} & the empty set & decided; and it bites three statements \\
\ref{sec:leg-emptyfamily} & the empty family is allowed & decided; the two formulations agree only under it \\
\ref{sec:leg-chain} & $\lll$ carries its chain & decided; \S5.5 is the sharpest instance \\
\ref{sec:leg-dimension} & ``dimension'' & \textbf{free} \\
\ref{sec:leg-linked} & ``linked instance'' & \textbf{free} \\
\ref{sec:leg-quotient} & quotients of subsets & decided; and see $\mathrm D17$ \\
\ref{sec:leg-notice} & sentences beginning ``Notice that'' & decided; promote to lemmas \\
\ref{sec:C10} & the dropped projective alternative & \textbf{a missing lemma}, proved \\
\hline
\end{tabular}
\end{center}

\subsection{The algebra $\mathbf Z_{p}$: not a convention, a theorem}
\label{sec:leg-Zp}

\texttt{main:1115} declares \emph{``In the paper every algebra $\mathbf Z_{p}$
belongs to $\Vn$ for a fixed $n$, hence the algebra $\mathbf Z_{p}$ is uniquely
defined''}, with $w^{\mathbf Z_{p}}=x_{1}+\dots+x_{n}\bmod p$. That operation is
idempotent only when $n\equiv1\pmod p$, so the declaration is a constraint on
$p$, and the source never says where the constraint comes from.

It comes from \emph{specialness}. If $\mathbf U$ is affine over $\mathbb Z_{p}$
with $|U|>1$ and $w$ is idempotent, weak near-unanimity and special of arity
$n$, then $w$ acts as $\sum_{i}c_{i}x_{i}+d$; idempotence gives
$\sum_{i}c_{i}=1$ and $d=0$; the WNU identities force all $c_{i}$ equal to a
single $c$ with $nc=1$; and specialness compares the coefficient of $y$ on the
two sides of $w(x,\dots,x,y)=w(x,\dots,x,w(x,\dots,x,y))$ to give $c^{2}=c$, so
$c\in\{0,1\}$, and $c=0$ contradicts $nc=1$. Hence $c=1$, $n\equiv1\pmod p$,
and $w$ acts as the sum.

\emph{Rendering.} Keep the two roles of $\mathbf Z_{p}$ apart: the abelian
\emph{group} $(\mathbb Z_{p};+,-)$, which supplies the relation
$x_{1}-x_{2}=x_{3}-x_{4}$ in the definition of a linear congruence and carries
no arity constraint; and the \emph{algebra}
$\mathbf Z_{p}=(\mathbb Z_{p};x_{1}+\dots+x_{n})\in\Vn$, which does. Only the
second occurs at \texttt{main:1300} ($\zeta\le\mathbf A\times\mathbf A\times
\mathbf Z_{p}$, a subalgebra of a product, which needs a common signature) and
at \texttt{main:3484}. Proof paper: \pkey{prop:p-divides},
\pkey{cor:Zp-in-Vn}, \pkey{haz:Zp}.

The proposition is not decoration: it is needed twice, and in both places its
absence leaves a statement that does not typecheck --- it is what makes the
definition of a perfect linear congruence well formed, and what upgrades the
group isomorphism of $\mathrm D16$ to an isomorphism of algebras.

\subsection{$\cover\sigma$ is a tolerance}\label{sec:leg-cover}

$\cover\sigma$ is defined (\texttt{main:1246}) as the minimal
$\delta\le\mathbf A\times\mathbf A$ with $\delta\supsetneq\sigma$ and $\delta$
stable under $\sigma$. Nothing in that makes it transitive.

\emph{Rendering.} $\cover\sigma$ is a reflexive symmetric subalgebra of
$\mathbf A^{2}$ and \emph{not} a congruence. Three facts need one line each and
none is in the source: the family of such $\delta$ is closed under
intersection unless the intersection is $\sigma$, which is exactly what
irreducibility forbids, so a least element exists; $(\cover\sigma)^{-1}$ is
another such relation, so symmetry follows by minimality --- used silently at
\texttt{SS:1538} (``switching $x_{1}$ and $x_{2}$''); and reflexivity is
$\sigma\subseteq\cover\sigma$.

\emph{The internal evidence.} That $\cover\sigma$ is a congruence is item (2)
of the definition of a linear congruence (\texttt{main:1371}) and conclusion
(1) of \texttt{LEMNontrivialReflexiveBridgeImplies} --- i.e.\ it is a property
of linear congruences, proved, not a property of irreducible ones. Typing it as
a congruence would make item (2) vacuous and collapse the linear/PC
distinction. Proof paper: \pkey{prop:cover}, \pkey{rem:cover-not-cong},
\pkey{haz:irreducible}.

\subsection{The empty set}\label{sec:leg-empty}

Read syntactically, $\varnothing$ is a subuniverse of every idempotent algebra,
is vacuously absorbing ($t(\varnothing,\dots,A,\dots,\varnothing)\subseteq
\varnothing$ holds because the left side is empty) and vacuously central (the
clause quantifies over $a\in A\setminus C$ against $\Sg(\varnothing)=
\varnothing$).

\emph{Rendering.} Follow \texttt{main:1629} --- \emph{``we do not allow empty
subuniverses''} --- and make $\lll$, $\subt{T}$, $\subte{T}$ relations between
\emph{nonempty} sets, with dotted variants for the empty case. Three statements
must then be restated, because they are true in the source only under the
vacuous reading; they are the first three rows of \S\ref{sec:restated}.

\subsection{The empty family is allowed}\label{sec:leg-emptyfamily}

Is $\sigma=A^{2}$ irreducible? The source gives two formulations of
irreducibility (\texttt{main:1237}), and \emph{they agree only under the
empty-family reading}, which settles the question.

\begin{itemize}\tightlist
\item Formulation 1: $\sigma$ is not an intersection of \emph{other} stable
  subalgebras of $\mathbf A^{2}$. The empty intersection is $A^{2}$, so
  $\sigma=A^{2}$ is reducible.
\item Formulation 2: there are no $S_{1},\dots,S_{k}\le\mathbf A/\sigma\times
  \mathbf A/\sigma$ with $0_{\mathbf A/\sigma}=\bigcap_{i}S_{i}$ and
  $S_{i}\neq0_{\mathbf A/\sigma}$. With $k=0$ this says
  $0_{\mathbf A/\sigma}\neq(A/\sigma)^{2}$, i.e.\ $\sigma\neq A^{2}$.
\end{itemize}

Forbid the empty family and formulation~2 no longer excludes $\sigma=A^{2}$
while formulation~1 still does.

\emph{Rendering.} Allow $k=0$. Then an irreducible congruence is automatically
proper, $\cover\sigma$ exists, and \cite{ZhukJACM}'s word ``proper'' --- which
\cite{ZhukSimplified} drops --- is recovered as a consequence rather than
reinstated by hand. Proof paper: \pkey{def:irreducible},
\pkey{prop:cover}(a).

\subsection{$\lll$ carries its chain}\label{sec:leg-chain}

$C\lll^{A}B$ is defined (\texttt{main:1588}) as the \emph{existence} of a chain
$C=B_{m}\subt{T_{m}}\dots\subt{T_{1}}B_{0}=B$; several \S5 proofs then speak of
\emph{``the dividing congruences coming from $C\lll^{A}B$''}
(\texttt{main:1600}), which is a function of the chain, not of the pair.

The sharpest instance is the $\TD\times\TD$ case of
\texttt{LEMTwoStableIntersection} (\texttt{SS:1840}). The step works only if
the chain witnessing $C_{2}\lll\mathbf A$ is the chain witnessing
$B_{2}\lll\mathbf A$ extended by the single step $C_{2}\subt{\TD(\sigma_{2})}
B_{2}$ --- that is what makes the dividing congruences of $C_{2}\lll\mathbf A$
equal to those of $B_{2}\lll\mathbf A$ together with $\sigma_{2}$, which is the
whole content of the appeal. The same case also applies the clause to $C_{2}$
rather than to $B_{2}$; against $B_{2}$ it does not follow, since both
alternatives are consistent with the hypotheses and the ``moreover'' clause
then says nothing.

\emph{Rendering.} $\lll$ carries its witness: a list of pairs (subset, typed
step with its congruence). Statements that mention ``the congruences coming
from $C\lll B$'' take the chain as an argument, extension of a chain by one
step is an operation on the data, and the propositional $\lll$ is its
truncation, used only where no congruence is named. Proof paper:
\pkey{def:lll}, \pkey{haz:lll-chain}.

\subsection{``Dimension'': free}\label{sec:leg-dimension}

$\prod_{i}\mathbb Z_{q_{i}}$ with distinct primes is not a vector space over
anything, so it has no dimension. Use the composition length of the finite
abelian group, which for $\mathbb Z_{p}^{n}$ is $n$ and agrees with the source
wherever the source is meaningful, and prove additivity along short exact
sequences. Nothing in \S5 uses the notion; it is confined to the codimension-one
argument. \textbf{The legislated reading is now unnecessary there:} the repaired
Step~2 never counts a dimension. It shows $L$ is the graph of an affine map and
reads $\Delta$ off as a coset of $\ker\lambda$ for a surjection $\lambda$ onto
$\mathbf Z_{p}$, and $\lambda$ kills every factor of order coprime to $p$ for
order reasons alone. ``Codimension one'' survives only as ``$\ker\lambda$ has
index $p$''. The notion is therefore retained for the source's benefit and
consumed nowhere. Proof paper: \pkey{lem:subuniverse-coset},
\pkey{haz:step2-linear-algebra}.

\subsection{``Linked instance'': free}\label{sec:leg-linked}

Two inequivalent definitions coexist: global connectivity of the value graph,
in the informal claim, and the per-variable condition in \S3.1. The second does
not imply the first, because a fragmented instance is per-variable linked in
each component.

\emph{Rendering.} Use the \S3.1 definition and carry non-fragmentation as a
separate hypothesis wherever the informal reading was doing that work. This is
not cosmetic: the retracted recursion-depth claim (\S\ref{sec:refuted-depth})
turned on exactly this pair of conditions at the \textsc{SolveNonlinked} call
site. Proof paper: \pkey{def:instance-linked}, \pkey{haz:linked}.

\subsection{Quotients of subsets}\label{sec:leg-quotient}

$B/\sigma$ denotes the \emph{image} $\{b/\sigma:b\in B\}$ in $\mathbf A/\sigma$,
for $B$ any subset of $A$ and $\sigma$ a congruence on $\mathbf A$ --- not the
quotient of $B$ by a congruence of $\mathbf B$. Consequently
$(C_{1}\cap\dots\cap C_{t})/\delta\subseteq\bigcap_{i}C_{i}/\delta$ always,
with equality \emph{false} in general.

\emph{Rendering.} Define the image operation once, prove only the inclusion,
and flag every site that needs the reverse. The reverse inclusion is
established on the fly inside clause (fm) of Propagation --- correctly, from
the saturation identity $\delta\cap B^{2}\subseteq\sigma_{i}\cap B^{2}$ --- and
must be extracted as a standalone lemma with its hypothesis visible. Proof
paper: \pkey{haz:quotient-rel}, and the (fm) clause of
\pkey{lem:propagation}.

This is the item the invalid repair of $\mathrm D17$ walked into
(\S\ref{sec:D17}). It is worth restating what that costs: cataloguing a trap is
not the same as being immune to it, and the only defence that worked was
someone else reading the argument.

\subsection{Sentences beginning ``Notice that''}\label{sec:leg-notice}

Four items, three of them from the source's own ``Notice that'' sentences, all
used later as lemmas.

\begin{enumerate}\alphlabels\tightlist
\item \texttt{main:1378}: the relation $S$ from item (3) of the definition of a
  linear congruence is a bridge from $\sigma$ to $\sigma$ with
  $\bcol S=\proj_{1,2}(S)=\proj_{3,4}(S)=\cover\sigma$. \textbf{Checked,
  true.} Clause (4) is $x_{1}-x_{2}=0\iff x_{3}-x_{4}=0$ inside each block;
  clause (3) holds because some block has $m\ge1$, else $\cover\sigma=\sigma$;
  and $\bcol S=\cover\sigma$ because $x-x=y-y$ always. Consumed at
  \texttt{SS:1360}. Proof paper: \pkey{lem:linear-bridge}.
\item \texttt{main:1386}: $\sigma$ is irreducible, PC or linear iff
  $0_{\mathbf A/\sigma}$ is. Needed to justify the ``factorize by $\sigma$ and
  assume $\sigma=0$'' opening used in three \S5.4 proofs. Proof paper:
  \pkey{lem:irred-quotient}.
\item \texttt{main:1550}: the two formulations of $\TS$-freeness agree.
  Trivial on unfolding $\subt{\TS}$ (\texttt{main:1522}), but it must be
  written, since $\TS$-freeness is a hypothesis of clause (m) of propagation
  modulo a congruence. Proof paper: \pkey{lem:sfree-equiv}.
\item Two assertions about rectangularity belong to the same list: that the
  parallelogram property implies rectangularity (``as it follows from the
  definitions''), and that the rectangular closure exists (asserted by the
  definite article in ``\emph{the} minimal rectangular relation''). Note also
  that ``$R$ is rectangular'' unqualified --- the form used in ``rectangular
  closure'' --- is never defined in the source; it means that every variable is
  rectangular. Proof paper: \pkey{lem:rect-basics}.
\end{enumerate}

\subsection{What legislating cost}\label{sec:leg-cost}

Six of the ten (\S\S\ref{sec:leg-cover}, \ref{sec:leg-empty},
\ref{sec:leg-emptyfamily}, \ref{sec:leg-chain}, \ref{sec:leg-quotient},
\ref{sec:leg-notice}) are decided by internal evidence. Two
(\S\S\ref{sec:leg-dimension}, \ref{sec:leg-linked}) are free choices where the
source is simply ambiguous and either reading can be made to work. Two
(\S\ref{sec:leg-Zp}, \S\ref{sec:C10}) are not conventions at all and have been
proved.
